\begin{center}
\textsc{Gauss} an \textsc{Olbers}.\quad Göttingen, 2.~September 1837.
\end{center}

\par\noindent\hspace*{\parindent}\mbox{.\ .\ .\ .\ .\ .}\quad \textsc{Gerling}, der seine trigonometrischen Messungen in Hessen jetzt abgeschlossen hat, hat jetzt noch eine Operation angesetzt, die zur Bestimmung des Längenunterschieds zwischen Göttingen und Mannheim dienen soll.
Auf zwei Bergen, dem Meisner und dem Feldberg, werden Signale gegeben. Die vom ersten Berg sind in meiner Sternwarte sichtbar, die vom zweiten in der Mannheimer Sternwarte, beide zugleich aber auf einem dazwischenliegenden Berg bei Marburg, wo \textsc{Gerling} mit einem Chronometer von \textsc{Kessel} beobachtet.
Die Signale sind nachts Pulverzeichen und
\par\smallskip\noindent\hfill{\footnotesize 48\textsuperscript{*}}\hspace*{.12\linewidth}\par
\clearpage
\noindent\makebox[0pt][l]{\fontsize{13}{16}\selectfont \textbf{380}}\hfill{\small\textbf{BRIEFWECHSEL MIT GERLING.}}\hfill\null\par\medskip
\noindent tagsüber heliotropische Zeichen.
Das Wetter ist nicht günstig. Hier sind bisher zwar schon recht viele Zeichen beider Arten beobachtet worden, doch \textsc{Gerling} hatte seinem letzten Brief zufolge noch fast nichts vom Feldberg her gesehen.\ \mbox{.\ .\ .\ .\ .\ .}

\begin{center}\rule{.24\linewidth}{.4pt}\end{center}
